\subsection{平成23年度 問6}
\label{q:h23-reg-applied-06}

\noindent\textbf{原典：} sources/H23/17\_applied\_science.pdf、PDF物理3頁。

\textbf{公式正答：} 3（single）。

\par\textit{原典本文の校正と独自解説は未作成。分類のみの登録項目。}

\subsection{平成25年度 問4}
\label{q:h25-reg-applied-04}

\noindent\textbf{原典：} sources/H25/17\_applied\_science.pdf、PDF物理3頁。

\textbf{公式正答：} 5（single）。

\par\textit{原典本文の校正と独自解説は未作成。分類のみの登録項目。}

\subsection{平成27年度 問4}
\label{q:h27-reg-applied-04}

\noindent\textbf{原典：} sources/H27/17\_applied\_science.pdf、PDF物理2頁。

\textbf{公式正答：} 4（single）。

\par\textit{原典本文の校正と独自解説は未作成。分類のみの登録項目。}

\subsection{平成28年度 問6}
\label{q:h28-reg-applied-06}

\noindent\textbf{原典：} sources/H28/17\_applied\_science.pdf、PDF物理2頁。

\textbf{公式正答：} 4（single）。

\par\textit{原典本文の校正と独自解説は未作成。分類のみの登録項目。}

\subsection{平成29年度 問5}
\label{q:h29-reg-applied-05}

\noindent\textbf{原典：} sources/H29/17\_applied\_science.pdf、PDF物理2頁。

\textbf{公式正答：} 4（single）。

\par\textit{原典本文の校正と独自解説は未作成。分類のみの登録項目。}

\subsection{令和元年度再試験 問1}
\label{q:r01-ret-applied-01}

\noindent\textbf{原典：} sources/R01\_retest/17\_applied\_science.pdf、PDF物理1頁。

\textbf{公式正答：} 1（single）。

\par\textit{原典本文の校正と独自解説は未作成。分類のみの登録項目。}

\subsection{令和2年度 問9}
\label{q:r02-reg-applied-09}

\noindent\textbf{原典：} sources/R02/17\_applied\_science.pdf、PDF物理3頁。

\textbf{公式正答：} 2（single）。

\par\textit{原典本文の校正と独自解説は未作成。分類のみの登録項目。}

\subsection{令和3年度 問9}
\label{q:r03-reg-applied-09}

\noindent\textbf{原典：} sources/R03/17\_applied\_science.pdf、PDF物理4頁。

\textbf{公式正答：} 1（single）。

\par\textit{原典本文の校正と独自解説は未作成。分類のみの登録項目。}

\subsection{令和5年度 問13}
\label{q:r05-reg-applied-13}

\noindent\textbf{原典：} sources/R05/17\_applied\_science.pdf、PDF物理5頁。

\textbf{公式正答：} 5（single）。

\par\textit{原典本文の校正と独自解説は未作成。分類のみの登録項目。}

\subsection{令和6年度 問16}
\label{q:r06-reg-applied-16}

\noindent\textbf{原典：} sources/R06/17\_applied\_science.pdf、PDF物理5頁。

\textbf{公式正答：} 3（single）。

\par\textit{原典本文の校正と独自解説は未作成。分類のみの登録項目。}

\subsection{令和7年度 問9}
\label{q:r07-reg-applied-09}

\noindent\textbf{原典：} sources/R07/17\_applied\_science.pdf、PDF物理3頁。

\textbf{公式正答：} 3（single）。

\par\textit{原典本文の校正と独自解説は未作成。分類のみの登録項目。}
